\section{Preliminaries}

\subsection{Notation}

Let $x = (x_1, \ldots, x_n)$ denote an input prompt and
$y = (y_1, \ldots, y_m)$ a response.
We write $y_{<t} \triangleq (y_1, \ldots, y_{t-1})$ for the
prefix up to step $t$.
We consider two LLMs: a student $\pi_\theta$
and a teacher $\pi_T$, each defining a next-token distribution
$\pi(\cdot \mid x, y_{<t})$ over a vocabulary $\mathcal{V}$.
We write $y \sim \pi_\theta(\cdot \mid x)$ for a response
sampled autoregressively from the student.
$\mathcal{D} = \{(x^{(i)}, y^{(i)})\}_{i=1}^N$ denotes a
fixed dataset of prompt–response pairs with teacher-generated
outputs, and
$\mathcal{D}_x \triangleq \{x^{(i)}\}_{i=1}^N$ the
corresponding prompt set. Knowledge distillation (KD) transfers knowledge from $\pi_T$ to $\pi_\theta$ 
by minimizing the divergence between the two distributions. A standard choice is 
the Kullback-Leibler (KL) divergence, defined for two distributions $P$ and 
$Q$ over $\mathcal{V}$ as $D_{\mathrm{KL}}(P \| Q) 
  = \sum_{v \in \mathcal{V}} P(v) \log \frac{P(v)}{Q(v)}$.





\subsection{On-Policy Distillation}
\label{sec:on_policy_distillation}

On-Policy Distillation (OPD) computes supervision on trajectories
sampled from the current student $\pi_\theta$. Given a prompt
$x \sim \mathcal{D}_x$, the student samples a response
$\hat{y} = (\hat{y}_1, \ldots, \hat{y}_{T})
  \sim \pi_\theta(\cdot \mid x)$,
where $T \triangleq |\hat{y}|$ denotes the rollout length.
Both models are then evaluated on the student-generated prefixes
$\hat{y}_{<t}$, yielding two next-token distributions at each
step~$t$:
$p_t(v) \triangleq \pi_\theta(v \mid x, \hat{y}_{<t})$ and
$q_t(v) \triangleq \pi_T(v \mid x, \hat{y}_{<t})$
for $v \in \mathcal{V}$.


A standard formulation minimizes the sequence-level reverse KL over
student-generated trajectories:
\begin{equation}
  \mathcal{L}_{\mathrm{OPD}}(\theta)
  =
  \mathbb{E}_{x \sim \mathcal{D}_x}
  \Bigl[
    D_{\mathrm{KL}}\bigl(\pi_\theta(\cdot \mid x)\;\|\;\pi_T(\cdot \mid x)\bigr)
  \Bigr].
  \label{eq:opd_seq_rkl_prelim}
\end{equation}

Using the autoregressive factorization, this sequence-level objective admits the exact token-level decomposition:
\begin{equation}
  \mathcal{L}_{\mathrm{OPD}}(\theta)
  =
  \mathbb{E}_{x \sim \mathcal{D}_x,\;
    \hat{y} \sim \pi_\theta(\cdot \mid x)}
  \left[
    \sum_{t=1}^{T}
      D_{\mathrm{KL}}(p_t \| q_t)
  \right].
  \label{eq:opd_exact_token_decomp}
\end{equation}
In practice, different implementations vary
in how this exact per-token reverse KL is computed: full-vocabulary OPD optimizes Eq.~\eqref{eq:opd_exact_token_decomp} directly, sampled-token OPD uses an unbiased Monte Carlo estimator of each token-level KL term, and top-$k$ OPD replaces the full-vocabulary KL with a subset-based approximation.
We now describe these three common supervision granularities.




\paragraph{Sampled-Token OPD.}
The most lightweight variant evaluates only the token sampled by the student, and is also the most common implementation in prior on-policy distillation work~\citep{lu2025onpolicydistillation,xiao2026mimo,yang2026learning}. Given $\hat{y}_t \sim p_t$, the per-token loss is
$\ell_t^{\mathrm{sample}}
  \triangleq \log p_t(\hat{y}_t) - \log q_t(\hat{y}_t)$,
aggregated as:
\begin{equation}
  \mathcal{L}_{\mathrm{OPD}}^{\mathrm{sample}}(\theta)
  = \mathbb{E}_{x \sim \mathcal{D}_x,\;
      \hat{y} \sim \pi_\theta(\cdot \mid x)}
    \left[\sum_{t=1}^{T}
      \ell_t^{\mathrm{sample}}\right].
  \label{eq:opd_sampled_obj_prelim}
\end{equation}
Since
$\mathbb{E}_{\hat{y}_t \sim p_t}
  [\ell_t^{\mathrm{sample}}]
  = D_{\mathrm{KL}}(p_t \| q_t)$,
each $\ell_t^{\mathrm{sample}}$ is an unbiased single-sample
estimator of the token-level reverse KL.


\paragraph{Full-Vocabulary OPD.}
At the other extreme, one computes the divergence over the entire
vocabulary at each prefix:
\begin{equation}
  \mathcal{L}_{\mathrm{OPD}}^{\mathrm{full}}(\theta)
  = \mathbb{E}_{x \sim \mathcal{D}_x,\;
      \hat{y} \sim \pi_\theta(\cdot \mid x)}
    \left[\sum_{t=1}^{T}
      D_{\mathrm{KL}}(p_t \| q_t)\right].
  \label{eq:opd_full_obj_prelim}
\end{equation}
This yields denser gradients compared to sampled-token OPD, at the cost of $O(BTM)$ memory for batch size $B$, sequence length $T$ and vocabulary size $M = |\mathcal{V}|$.


\paragraph{Top-$k$ OPD.}
Top-$k$ OPD provides an intermediate design between sampled-token and
full-vocabulary OPD by restricting the divergence computation to a subset
$S_t \subseteq \mathcal{V}$. Here we focus on the student top-$k$ variant,
which selects the $k$ tokens assigned the highest probability under the
student, namely $S_t = \operatorname{TopK}(p_t, k)$.
Define the renormalized student and teacher distributions on $S_t$ as:
\[
\bar{p}_t^{(S_t)}(v)
=
\frac{p_t(v)\,\mathbf{1}[v \in S_t]}
     {\sum_{u \in S_t} p_t(u)},
\qquad
\bar{q}_t^{(S_t)}(v)
=
\frac{q_t(v)\,\mathbf{1}[v \in S_t]}
     {\sum_{u \in S_t} q_t(u)}.
\]
Distillation is then performed by minimizing the subset KL divergence $ D_{\mathrm{KL}}\!\bigl(\bar{p}_t^{(S_t)} \,\|\, \bar{q}_t^{(S_t)}\bigr)$, yielding the trajectory-level objective:
\begin{equation}
  \mathcal{L}_{\mathrm{OPD}}^{\mathrm{top\text{-}k}}(\theta)
  =
  \mathbb{E}_{x \sim \mathcal{D}_x,\;
      \hat{y} \sim \pi_\theta(\cdot \mid x)}
  \left[
    \sum_{t=1}^{T}
    D_{\mathrm{KL}}\!\bigl(\bar{p}_t^{(S_t)} \,\|\, \bar{q}_t^{(S_t)}\bigr)
  \right].
  \label{eq:opd_topk_obj_prelim}
\end{equation}

This formulation discards mass outside $S_t$ and therefore remains an approximation to the full-vocabulary reverse KL, but it substantially reduces teacher-query cost while preserving multi-token supervision on the student's high-probability region.

\subsection{Dynamic Metrics}
\label{sec:dynamic_metrics}

We define the student's and teacher's top-$k$ sets at step $t$
as $S_t^{(p)} = \operatorname{TopK}(p_t, k)$ and
$S_t^{(q)} = \operatorname{TopK}(q_t, k)$, respectively.
The following metrics are monitored throughout OPD training in later experiments.


\paragraph{Overlap Ratio.}
This metric quantifies the alignment between the student's and teacher's candidate spaces. It is defined as the average proportion of tokens that appear simultaneously in both the student's and the teacher's top-$k$ sets:
\begin{equation}
\mathcal{M}_{\text{overlap}}
\triangleq
\mathbb{E}_{t} \left[ \frac{|S_t^{(p)} \cap S_t^{(q)}|}{k} \right].
\label{eq:metric_overlap}
\end{equation}
A low overlap ratio indicates that the student's probability mass is concentrated on a disjoint set of tokens from the teacher, suggesting significant policy divergence or ``mode mismatch''. Conversely, a ratio nearing $1.0$ implies the student has successfully located the teacher's support region.

\paragraph{Overlap-Token Advantage.}
To measure distributional agreement within the overlap tokens,
we define
$A_t(v) \triangleq \bar{p}_t(v)
  (\log \bar{q}_t(v) - \log \bar{p}_t(v))$
where $\bar{p}_t, \bar{q}_t$ are the renormalized student and
teacher distributions over $S_t^{(p)} \cap S_t^{(q)}$.
The metric averages this quantity:
\begin{equation}
  \mathcal{M}_{\text{adv}}
  \triangleq \mathbb{E}_{t}\!\left[
      \frac{1}{|S_t^{(p)} \cap S_t^{(q)}|}
      \sum_{v \in S_t^{(p)} \cap S_t^{(q)}}
        A_t(v)
    \right].
  \label{eq:metric_adv}
\end{equation}
A value close to zero indicates high-quality alignment where the student places mass on teacher-preferred tokens with appropriate confidence. Conversely, a large negative value indicates that within the intersection, the student is overconfident compared to the teacher (high $p_t$ but lower $q_t$).


\paragraph{Entropy and Entropy Gap.}
To monitor the distributional properties of the policies, we track the entropy of both the student $H(p_t)$ and the teacher $H(q_t)$ on the student's rollouts, and define the entropy gap as:
\begin{equation}
\Delta H_t = \left| H(q_t) - H(p_t) \right|.
\label{eq:metric_entropy_gap}
\end{equation}
$\Delta H_t$ is a state-specific indicator of mode alignment.
A large gap suggests a substantial mismatch between the student and teacher in confidence and diversity over the same visited states, while convergence toward zero indicates that the student has matched the teacher's uncertainty profile along its generated trajectories.














































